\title{Group Sequence Policy Optimization}

\begin{abstract}
We present Group Sequence Policy Optimization (GSPO), a reinforcement learning algorithm specifically developed for large language models, emphasizing stability, computational efficiency, and strong performance. Departing from earlier methods that leverage token-level importance ratios, GSPO establishes its importance ratio over entire sequences and applies clipping, reward assignment, and optimization at the sequence level. Our results reveal that GSPO not only delivers better training efficiency and performance than the GRPO algorithm but also significantly enhances the stability of Mixture-of-Experts (MoE) RL training. Moreover, GSPO offers opportunities to streamline the architecture of RL systems. These advantages have played a key role in the notable advances seen in the latest Qwen3 models.
\end{abstract}

\section{Introduction}
\label{sec:intro}

Reinforcement learning (RL) has become a central approach for enhancing the scalability of language models \citep{o1,dpsk-r1,qwq32b,qwen3}. When applied at scale, RL equips these models with the ability to solve complex tasks, such as high-level mathematical problems and advanced programming challenges, by enabling more profound and extensive reasoning.

A fundamental requirement for advancing RL with increased computational resources is the preservation of steady and resilient training behavior. Yet, prevailing RL techniques, such as GRPO \citep{grpo}, suffer from pronounced instability during the training of very large language models, frequently leading to drastic and sometimes irrecoverable model degradation or collapse \citep{qwen3, minimax-m1}. These instability issues present a significant obstacle to further enhancing language model abilities via continued RL optimization.

In this study, we demonstrate that the core instability observed in GRPO arises from a fundamental flaw: the algorithm improperly applies importance sampling weights, resulting in their breakdown. This flaw introduces substantial variance into the training process, with the noise intensifying as response sequences lengthen. The effect is further compounded by the clipping procedure, which accelerates the trajectory toward model collapse.

To mitigate these foundational problems, we introduce \textbf{Group Sequence Policy Optimization (GSPO)}, a novel reinforcement learning approach tailored for training large language models. GSPO's central contribution is its principled and theoretically sound definition of the importance ratio, constructed from sequence likelihood \citep{click}, thus upholding the essential tenets of importance sampling. Furthermore, GSPO determines normalized rewards by evaluating the advantages across a set of responses to each query, thereby maintaining a consistent alignment between sequence-level reward allocation and the corresponding optimization process.

Our empirical results indicate that GSPO substantially outperforms GRPO across key metrics such as training stability, efficiency, and overall effectiveness. Notably, GSPO addresses the inherent stability issues associated with reinforcement learning on large Mixture-of-Experts (MoE) models without resorting to intricate stabilization methods, suggesting its promise in streamlining RL infrastructures. These strengths have played a vital role in achieving marked performance gains in the latest Qwen3 models. We anticipate that GSPO will serve as a reliable and scalable basis for future progress in large-scale RL training of language models.

\section{Preliminaries}

\paragraph{Notation}
Throughout this work, we represent an autoregressive language model with parameters $\theta$ as a policy $\pi_\theta$. Here, the symbol $x$ is used for a query, while $\mathcal{D}$ refers to the set of queries. For a response $y$ corresponding to query $x$, the probability assigned by the policy $\pi_\theta$ is written as $\pi_\theta (y | x)=\prod_{t=1}^{|y|} \pi_\theta (y_t | x, y_{<t} )$, where $|y|$ stands for the length of $y$ in tokens. The quality or correctness of a query-response tuple $(x, y)$ can be assessed by a verifier $r$, generating a reward value $r(x, y) \in [0, 1]$.

\paragraph{Proximal Policy Optimization (PPO)} Using samples generated from the old policy $\pi_{\theta_\text{old}}$, PPO \citep{ppo} constrains the policy update within a proximal region of the old policy through the clipping mechanism. Specifically, PPO employs the following objective for policy optimization (we omit the KL regularization term hereinafter for brevity, as it is not the focus of this paper):

\begin{align}
\mathcal{J}_\text{PPO}(\theta) = \mathbb{E}_{ x \sim \mathcal{D},\, y \sim \pi_{\theta_\text{old}}( \cdot | x) }
\left[ \frac{1}{|y|} \sum_{t=1}^{|y|} 
\min \left( w_{t}(\theta) \widehat{A}_{t},  \, \mathrm{clip} \left( w_{t}(\theta), 1 - {\varepsilon}, 1 + {\varepsilon}\right) \widehat{A}_{t} \right)
\right],
\end{align}

Here, the token importance ratio $w_{t}(\theta)$ is calculated as $
w_{t}(\theta) = \frac{ \pi_{\theta} (y_{t} | x, y_{<t}) }{ \pi_{\theta_\text{old}} (y_{t} | x,y_{<t})} $, with the advantage $\widehat{A}_{t}$ for $y_t$ being provided by a separate value model. The variable $\varepsilon$ serves as the threshold for clipping the importance ratios.

A fundamental obstacle faced by PPO in real-world applications is its significant dependence on the value model. Typically, this value model matches the policy model in size, resulting in substantial demands on both memory and computation. Additionally, the overall success of the algorithm is closely tied to the accuracy of the value predictions. Not only is developing a dependable value model a difficult task in itself, but achieving its effectiveness for longer outputs and more sophisticated problem domains intensifies this difficulty.

\paragraph{Group Relative Policy Optimization (GRPO)} GRPO \citep{grpo} bypasses the need for the value model by computing the relative advantage of each response within a group of responses to the same query. Specifically, GRPO optimizes the following objective:

\begin{align}
\label{equ:grpo}
\mathcal{J}_\text{GRPO}(\theta) = \mathbb{E}_{ x \sim \mathcal{D},\, \{y_i\}_{i=1}^G \sim \pi_{\theta_\text{old}}( \cdot | x) }
\left[ \frac{1}{G} \sum_{i=1}^{G} \frac{1}{|y_i|} \sum_{t=1}^{|y_i|} 
\min \left( w_{i,t}(\theta) \widehat{A}_{i,t},  \, \mathrm{clip} \left( w_{i,t}(\theta), 1 - {\varepsilon}, 1 + {\varepsilon}\right) \widehat{A}_{i,t} \right)
\right],
\end{align}

where $G$ is the number of generated responses to each query $x$ (i.e., the group size), and the importance ratio $w_{i,t}(\theta)$ and advantage $\widehat{A}_{i,t}$ of token $y_{i,t}$ are:

\begin{align}
    w_{i,t}(\theta)=\frac{ \pi_{\theta} (y_{i,t} | x, y_{i,<t}) }{ \pi_{\theta_\text{old}} (y_{i,t} | x,y_{i,<t})},\quad\ 
    \widehat{A}_{i,t} = \widehat{A}_{i} = \frac{r(x, y_i) - \mathrm{mean} \left( \{ r(x, y_i) \}_{i=1}^G \right) }{ \mathrm{std} \left( \{ r(x, y_i) \}_{i=1}^G \right) },
\end{align}

respectively, where all the tokens in $y_i$ share the same advantage as $\widehat{A}_{i}$.

\section{Motivation}
\label{sec:motivation}

The increase in model parameters, the adoption of sparsity strategies such as those in Mixture-of-Experts architectures, and the generation of longer responses all require deploying larger rollout batch sizes to fully utilize computational resources during RL. To enhance the efficiency with which samples contribute to learning, it is a common strategy to divide these substantial batches into several smaller mini-batches for the purpose of performing gradient updates. However, this mini-batching leads inherently to an off-policy learning framework, whereby the generated responses $y$ originate from an earlier version of the policy, $\pi_{\theta_\text{old}}$, instead of the present policy $\pi_\theta$ targeted for improvement. Consequently, mechanisms like clipping in PPO and GRPO are essential, as they limit the contribution of samples that diverge excessively from the current policy, thereby maintaining the integrity of gradient estimates.

While mechanisms like clipping aim to manage this off-policy discrepancy, we identify a more fundamental issue in GRPO: \textit{its objective is ill-posed}. This problem becomes particularly acute when training large models on long-response tasks, leading to catastrophic model collapse. The ill-posed nature of the GRPO objective stems from a misapplication of importance sampling weights. The principle of importance sampling is to estimate the expectation of a function $f$ under a target distribution $\pi_\text{tar}$ by re-weighting samples drawn from a behavior distribution $\pi_\text{beh}$:

\begin{align}
\mathbb{E}_{z \sim \pi_\text{tar}} \left[ f(z) \right] = \mathbb{E}_{z \sim \pi_\text{beh}} \left[ \frac{ \pi_\text{tar} (z) }{ \pi_\text{beh} (z)} \, f(z) \right].
\end{align}

A key aspect of this approach is the necessity to average across a substantial number of samples ($N \gg 1$) drawn from the behavior distribution $\pi_\text{beh}$ in order for the importance weight $\frac{ \pi_\text{tar} (z) }{ \pi_\text{beh} (z)}$ to adequately adjust for mismatches between distributions.

GRPO, by contrast, implements the importance weight $\frac{ \pi_{\theta} (y_{i,t} | x, y_{i,<t}) }{ \pi_{\theta_\text{old}} (y_{i,t} | x,y_{i,<t})}$ independently at each token position $t$. Since this weight utilizes only a single sample $y_{i,t}$ from the next-token distribution $\pi_{\theta_\text{old}}( \cdot | x, y_{i,<t})$, it is ineffective at correcting for distributional discrepancies as intended. Rather than facilitating proper adjustment, it injects considerable high-variance noise into the training gradients. Over lengthy sequences, this noise aggregates and is further intensified by the application of clipping strategies. In practice, we have seen this instability give rise to catastrophic model collapse. Once such collapse happens, remedial actions such as continuing training, restoring previous checkpoints, exhaustively tuning hyperparameters (e.g., clipping boundaries), prolonging generated sequence lengths, or modifying RL queries do not reverse the degeneration.

This finding reveals a deeper flaw in the design of GRPO. The ineffectiveness of the per-token importance weighting highlights a fundamental mismatch: \textbf{the granularity at which the optimization objective is defined should be consistent with the granularity of the reward signal}. Since rewards are allocated at the sequence level, conducting off-policy corrections at the token level is inherently fraught. This insight leads us to abandon the token-level strategy and instead pursue sequence-level importance weighting and optimization, aligning the level of off-policy correction with that of the reward structure.

\section{Algorithm}

\subsection{GSPO: Group Sequence Policy Optimization}

Although using the token-level importance weight $\frac{ \pi_{\theta} (y_{i,t} | x, y_{i,<t}) }{ \pi_{\theta_\text{old}} (y_{i,t} | x,y_{i,<t})}$ presents challenges within GRPO, we note that, for language generation tasks, the \textit{sequence-level} importance weight $\frac{ \pi_{\theta} (y | x) }{ \pi_{\theta_\text{old}} (y | x)}$ carries a precise theoretical significance. Specifically, it quantifies the degree of divergence between the generated response $y$—sampled according to $\pi_{\theta_\text{old}} (\cdot | x)$—and the current model distribution $\pi_{\theta} (\cdot | x)$. This measure is inherently aligned with rewards computed over whole sequences and offers a principled metric for implementing the clipping procedure.

Based on this straightforward observation, we propose the \textbf{Group Sequence Policy Optimization (GSPO)} algorithm. GSPO employs the following sequence-level optimization objective:

\begin{align}
\mathcal{J}_\text{GSPO} (\theta) =
\mathbb{E}_{ x \sim \mathcal{D},\, \{y_i\}_{i=1}^G \sim \pi_{\theta_\text{old}}( \cdot | x) }
\left[ 
\frac{1}{G} \sum_{i=1}^{G}
\min \left( s_{i}(\theta)  \widehat{A}_{i},  \, \mathrm{clip} \left( s_{i}(\theta), 1 - {\varepsilon}, 1 + {\varepsilon} \right) \widehat{A}_{i} \right) 
\right],
\label{equ:gspo}
\end{align}

where we adopt the group-based advantage estimation:

\begin{align}
\widehat{A}_{i} = \frac{r(x, y_i) - \mathrm{mean} \left( \{ r(x, y_i) \}_{i=1}^G \right) }{ \mathrm{std} \left( \{ r(x, y_i) \}_{i=1}^G \right) },
\end{align}

and define the importance ratio $s_{i}(\theta)$ based on sequence likelihood \citep{click}:

\begin{align}
s_{i}(\theta) = \left( \frac{ \pi_{\theta} (y_i | x) }{ \pi_{\theta_\text{old}} (y_i | x)} \right)^{\frac{1}{|y_i|}}
=
\exp \left( \frac{1}{|y_i|} \sum_{t=1}^{|y_i|} \log \frac{ \pi_{\theta} (y_{i,t} | x, y_{i,<t}) }{ \pi_{\theta_\text{old}} (y_{i,t} | x,y_{i,<t})} \right).
\end{align}

As a result, GSPO utilizes response-level clipping rather than token-level clipping, thereby filtering out excessively ``off-policy'' samples from the gradient computation and aligning both with the rewards given at the sequence level and the optimization objective. Importantly, we implement length normalization in $s_{i}(\theta)$, which serves to both decrease variance and confine $s_{i}(\theta)$ within a consistent numerical interval. Without this normalization, significant changes in the likelihood of just a few tokens could lead to large swings in the overall sequence-level importance ratio, necessitating different clipping bounds for responses of varying lengths. Additionally, we observe that the clipping ranges employed in GSPO are generally orders of magnitude apart from those in prior methods such as GRPO, a result of the fundamentally different definitions for importance ratios in each approach.

\subsection{Gradient Analysis}
\label{subsec:gradient}

We can derive the gradient of the GSPO objective as follows (clipping is omitted for brevity):

\begin{align}
\nabla_{\theta} \mathcal{J}_\text{GSPO} (\theta)
=&\ 
\nabla_{\theta} \mathbb{E}_{ x \sim \mathcal{D},\, \{y_i\}_{i=1}^G \sim \pi_{\theta_\text{old}}( \cdot | x) }
\left[ 
\frac{1}{G} \sum_{i=1}^{G} s_{i}(\theta) \widehat{A}_{i}
\right] \\
= &\ 
\mathbb{E}_{ x \sim \mathcal{D},\, \{y_i\}_{i=1}^G \sim \pi_{\theta_\text{old}}( \cdot | x) }
\left[ 
\frac{1}{G} \sum_{i=1}^{G}
s_{i}(\theta) \widehat{A}_{i}
\cdot \nabla_{\theta} \log s_{i}(\theta)
\right] \\
=&\ 
\mathbb{E}_{ x \sim \mathcal{D},\, \{y_i\}_{i=1}^G \sim \pi_{\theta_\text{old}}( \cdot | x) }
\left[ 
\frac{1}{G} \sum_{i=1}^{G} 
\left( \frac{ \pi_{\theta} (y_i | x) }{ \pi_{\theta_\text{old}} (y_i | x)} \right)^{\frac{1}{|y_i|}} \widehat{A}_{i} 
\cdot \frac{1}{|y_i|} \sum_{t=1}^{|y_i|} \nabla_{\theta} \log \pi_{\theta} (y_{i,t} | x, y_{i,<t}) 
\right].
\label{equ:gspo_grad}
\end{align}

For comparison, the gradient of the GRPO objective is as follows (note that $\widehat{A}_{i,t} = \widehat{A}_{i}$):

\begin{align}
\nabla_{\theta} \mathcal{J}_\text{GRPO} (\theta) 
=&\ 
\nabla_{\theta} \mathbb{E}_{ x \sim \mathcal{D},\, \{y_i\}_{i=1}^G \sim \pi_{\theta_\text{old}}( \cdot | x) }
\left[ \frac{1}{G} \sum_{i=1}^{G} \frac{1}{|y_i|} \sum_{t=1}^{|y_i|} 
w_{i,t}(\theta) \widehat{A}_{i,t}
\right] \\
=&\
\mathbb{E}_{ x \sim \mathcal{D},\, \{y_i\}_{i=1}^G \sim \pi_{\theta_\text{old}}( \cdot | x) }
\left[ \frac{1}{G} \sum_{i=1}^{G} \widehat{A}_{i} 
\cdot \frac{1}{|y_i|} \sum_{t=1}^{|y_i|} 
\frac{ \pi_{\theta} (y_{i,t} | x, y_{i,<t}) }{ \pi_{\theta_\text{old}} (y_{i,t} | x,y_{i,<t})} 
\nabla_{\theta} \log \pi_{\theta} (y_{i,t} | x, y_{i,<t})  
\right].
\end{align}

The key difference between GSPO and GRPO centers on \textit{the manner in which gradient contributions from token log likelihoods are weighted}. In the case of GRPO, each token’s contribution is adjusted by an ``importance weight'' factor given by $\frac{ \pi_{\theta} (y_{i,t} | x, y_{i,<t}) }{ \pi_{\theta_\text{old}} (y_{i,t} | x,y_{i,<t})}$. These variable weights, which fall within $(0, 1+\varepsilon]$ when $\widehat{A}_{i} >0$ or $[1-\varepsilon, +\infty)$ if $\widehat{A}_{i} < 0$, can differ significantly between tokens. As a result, their cumulative effect over training may result in instability and unforeseen dynamics. In contrast, GSPO assigns an equal weight to every token within a response, thereby removing the instability introduced by the importance weights in GRPO.

\subsection{GSPO-token: A Token-level Objective Variant}

In scenarios like multi-turn RL, we may desire a finer-grained advantage adjustment than the sequence level. To this end, we introduce a token-level objective variant of GSPO, namely \textbf{GSPO-token}, to allow token-wise advantage customization:

\begin{align}
\mathcal{J}_\text{GSPO-token}(\theta)
=
\mathbb{E}_{ x \sim \mathcal{D},\, \{y_i\}_{i=1}^G \sim \pi_{\theta_\text{old}}( \cdot | x) }
\left[ 
\frac{1}{G} \sum_{i=1}^{G} \frac{1}{|y_i|} \sum_{t=1}^{|y_i|} 
\min \left( s_{i,t}(\theta) \widehat{A}_{i,t},  \, \mathrm{clip} \left( s_{i,t}(\theta), 1 - {\varepsilon}, 1 + {\varepsilon} \right) \widehat{A}_{i,t} \right) 
\right],
\label{equ:gspo-token}
\end{align}

where

\begin{align}
s_{i,t}(\theta) 
= \mathrm{sg} \left[ s_{i}(\theta) \right]  \cdot \frac{ \pi_{\theta} (y_{i,t} | x, y_{i,<t}) }{ \mathrm{sg} \left[ \pi_{\theta} (y_{i,t} | x, y_{i,<t}) \right] },
\end{align}

and $\mathrm{sg}[\cdot]$ denotes only taking the numerical value but stopping the gradient, corresponding to the \texttt{detach} operation in PyTorch. The gradient of GSPO-token can be derived as:

\begin{align}
\nabla_{\theta} \mathcal{J}_\text{GSPO-token}(\theta)
=&\ 
\nabla_{\theta} \mathbb{E}_{ x \sim \mathcal{D},\, \{y_i\}_{i=1}^G \sim \pi_{\theta_\text{old}}( \cdot | x) }
\left[ 
\frac{1}{G} \sum_{i=1}^{G}
\frac{1}{|y_i|} \sum_{t=1}^{|y_i|} 
s_{i,t}(\theta) \widehat{A}_{i,t}
\right] \\
=&\ 
\mathbb{E}_{ x \sim \mathcal{D},\, \{y_i\}_{i=1}^G \sim \pi_{\theta_\text{old}}( \cdot | x) }
\left[ 
\frac{1}{G} \sum_{i=1}^{G} s_i (\theta) \cdot \frac{1}{|y_i|} \sum_{t=1}^{|y_i|} 
\widehat{A}_{i,t} \frac{ \nabla_{\theta} \pi_{\theta} (y_{i,t} | x, y_{i,<t}) }{ \pi_{\theta} (y_{i,t} | x, y_{i,<t}) }
\right] \\
=&\ 
\mathbb{E}_{ x \sim \mathcal{D},\, \{y_i\}_{i=1}^G \sim \pi_{\theta_\text{old}}( \cdot | x) }
\left[ 
\frac{1}{G} \sum_{i=1}^{G}  \left( \frac{ \pi_{\theta} (y_i | x) }{ \pi_{\theta_\text{old}} (y_i | x)} \right)^{\frac{1}{|y_i|}}
\cdot \frac{1}{|y_i|} \sum_{t=1}^{|y_i|}
\widehat{A}_{i,t} \nabla_{\theta} \log \pi_{\theta} (y_{i,t} | x, y_{i,<t})  
\right].
\label{equ:gspo-token_grad}
\end{align}

It is important to observe that the fraction $\frac{ \pi_{\theta} (y_{i,t} | x, y_{i,<t}) }{ \mathrm{sg} \left[ \pi_{\theta} (y_{i,t} | x, y_{i,<t}) \right] }$ always evaluates to 1, thereby ensuring that $s_{i,t}(\theta)$ and $s_{i}(\theta)$ are equal in terms of their computed values. A direct comparison between Equation~\eqref{equ:gspo} and~\eqref{equ:gspo-token}, as well as between Equation~\eqref{equ:gspo_grad} and~\eqref{equ:gspo-token_grad}, reveals that GSPO-token and GSPO coincide exactly in the optimization objective, clipping criteria, and theoretical gradient, provided that all token-level advantages within a given response $y_i$ are set identically (that is, $\widehat{A}_{i,t} = \widehat{A}_{i}$ for all $t$). Nevertheless, GSPO-token offers the enhanced capability to specify distinct advantage values on a per-token basis.

\section{Experiments and Discussion}

\subsection{Empirical Results}

We conduct experiments utilizing a cold-start model that has been fine-tuned from Qwen3-30B-A3B-Base, and present both training reward trajectories and performance trends across AIME'24 (with average Pass@1 over 32 samples), LiveCodeBench (202410-202502, measured as average Pass@1 over 8 samples), as well as CodeForces (Elo Rating) evaluation benchmarks. Throughout the RL optimization, each rollout batch is subdivided into four smaller mini-batches, which are then used for separate gradient updates. For GSPO, we configure the left and right clipping bounds in Equation~\eqref{equ:gspo} to 3e-4 and 4e-4, respectively. As a baseline, we use GRPO, tuning its left and right clipping thresholds in Equation~\eqref{equ:grpo} to 0.2 and 0.27, respectively, to guarantee the validity of our comparison. It should be emphasized that normal convergence for MoE RL under GRPO requires adoption of the Routing Replay training procedure—a detail elaborated in \S~\ref{subsec:routing_replay}—whereas \textit{GSPO eliminates the requirement for this additional strategy}.

Figure~\ref{fig:results} illustrates that GSPO enables consistently stable training dynamics. Notably, \textbf{GSPO facilitates ongoing gains in performance by scaling up training compute, periodically refreshing the query set, and increasing the generation length}. In addition, GSPO exhibits greater efficiency in training when compared to GRPO, attaining higher training accuracy and superior benchmark outcomes under equivalent compute resources and query consumption. Lastly, the RL training of the most recent Qwen3 models using GSPO demonstrates its effectiveness, offering robust evidence of GSPO’s capacity to harness the benefits of RL scaling for large language models.

\subsection{Curious Observation on Clipping Fractions}

An important difference between GSPO and GRPO lies in GSPO's approach of clipping at the response level rather than at the token level. As depicted in Figure~\ref{fig:clipping}, there is a disparity of nearly two orders of magnitude in the proportion of clipped tokens between GSPO and GRPO, a gap that persists even when modifying the clipping thresholds. Nonetheless, even though GSPO excludes far more tokens from training due to extensive clipping, it surpasses GRPO in terms of training efficiency. This seemingly paradoxical result—where removing a larger share of tokens actually yields better training outcomes—suggests that GRPO’s per-token gradient estimates are fundamentally more variable and less effective for exploiting available samples. On the other hand, GSPO’s method of applying clipping to entire sequences offers a more stable and productive learning signal.

\subsection{Benefit of GSPO for MoE Training}
\label{subsec:routing_replay}

\paragraph{Background}
In contrast with RL training for dense architectures, the sparse expert selection inherent in MoE models brings about distinct challenges concerning training stability. Notably, our observations indicate that when utilizing the GRPO algorithm, the \textit{expert-activation volatility} present in MoE models may obstruct proper RL convergence. More precisely, after applying one or more rounds of gradient updates, the ensemble of experts triggered to produce the same output can shift considerably. For instance, in experiments with the 48-layer Qwen3-30B-A3B-Base model, approximately 10\% of the experts chosen for the same rollout input under the new policy $\pi_{\theta}$ differ from those selected under the previous policy $\pi_{\theta_\text{old}}$ following each RL gradient update. This volatility intensifies in deeper MoE structures, resulting in significant fluctuations—and potential invalidation—of the token-level importance ratios $w_{i,t}(\theta) = \frac{ \pi_{\theta} (y_{i,t} | x, y_{i,<t}) }{ \pi_{\theta_\text{old}} (y_{i,t} | x,y_{i,<t})}$, as elaborated in \S~\ref{sec:motivation} and~\ref{subsec:gradient}. Such instability, in turn, impedes the standard progress of RL training convergence.

\paragraph{Our Previous Approach} In our earlier work addressing this problem, we utilized a training method called \textbf{Routing Replay}. This approach involves recording the set of experts activated under $\pi_{\theta_\text{old}}$, and subsequently forcing $\pi_{\theta}$ to employ these same routing patterns while calculating the importance ratios $w_{i,t}(\theta) = \frac{ \pi_{\theta} (y_{i,t} | x, y_{i,<t}) }{ \pi_{\theta_\text{old}} (y_{i,t} | x,y_{i,<t})}$. As a result, both $\pi_{\theta} (y_{i,t} | x, y_{i,<t})$ and $\pi_{\theta_\text{old}} (y_{i,t} | x, y_{i,<t})$ evaluate each token $y_{i,t}$ using identical expert subnetworks. This preserves the stability of the token-wise importance ratios and guarantees that gradient updates are performed consistently on the same activated subnetwork. As depicted in Figure~\ref{fig:routing_replay}, Routing Replay proves to be a crucial mechanism for achieving stable convergence during GRPO training for MoE models.

\paragraph{Benefit of GSPO} While Routing Replay allows for successful GRPO training of MoE models, its strategy of repeatedly using prior routing decisions adds extra memory and communication costs, and may constrain the effective expressiveness of the MoE architecture. In comparison, as depicted in Figure~\ref{fig:results}, GSPO removes the reliance on Routing Replay and straightforwardly calculates the importance ratios $s_i(\theta)$, achieving conventional convergence and robust optimization. The central observation is that GSPO restricts its attention to the sequence likelihood, $\pi_{\theta} (y_i | x)$, rather than to the likelihood of individual tokens, $\pi_{\theta} (y_{i,t} | x, y_{i,<t})$, making it less vulnerable to token-level variations. Given that the MoE model consistently retains its language modeling strengths, the value of the sequence likelihood remains relatively stable. Overall, GSPO addresses the volatility in expert activation intrinsic to MoE architectures at a fundamental level, thereby negating the need for auxiliary mechanisms like Routing Replay. This leads to a more streamlined and robust training procedure, empowering the model to utilize its full representational potential without imposed restrictions.

\subsection{Benefit of GSPO for RL Infrastructure}

Due to inconsistencies in numerical precision between training engines (such as Megatron) and inference engines (such as SGLang and vLLM), the standard approach in practice involves recalculating the likelihoods of sampled responses for the prior policy $\pi_{\theta_\text{old}}$ using the training engine. In contrast, GSPO relies on sequence-level likelihoods instead of token-level ones for its optimization objective, and as a result, sequence-level metrics are inherently more resilient to precision mismatches. This property allows GSPO to make direct use of likelihoods produced by the inference engine for optimization, eliminating the need to recompute them with the training engine. Such a feature is particularly advantageous in use cases like partial rollout, multi-turn reinforcement learning, and frameworks where training and inference are separated.

\section{Conclusion}

We introduce Group Sequence Policy Optimization (GSPO), an innovative reinforcement learning method tailored for the training of large language models. Adhering to the foundational concepts of importance sampling, GSPO establishes importance ratios grounded in the likelihood of entire sequences and implements sequence-level mechanisms for clipping, reward assignment, and optimization. In comparison to GRPO, GSPO delivers considerably enhanced training stability, greater efficiency, and improved outcomes, and proves especially effective for conducting large-scale reinforcement learning with MoE models. This method underpins the remarkable advancements observed in the newest Qwen3 models. Leveraging the scalability of GSPO, our research will further scale reinforcement learning approaches, aiming to catalyze substantial progress in artificial intelligence.

\end{document}